\subsection{A General Coordinate Formula in the Plane}

\begin{derivationbox}[Eliminating the basis coefficients]
Starting from
\[
\begin{aligned}
 a_1c_1+b_1c_2&=x,\\
 a_2c_1+b_2c_2&=y,
\end{aligned}
\]
multiply the first equation by $b_2$ and the second by $b_1$:
\[
\begin{aligned}
 a_1b_2c_1+b_1b_2c_2&=xb_2,\\
 a_2b_1c_1+b_1b_2c_2&=yb_1.
\end{aligned}
\]
Subtracting eliminates $c_2$:
\[
(a_1b_2-a_2b_1)c_1=xb_2-yb_1.
\]
Similarly, eliminating $c_1$ gives
\[
(a_1b_2-a_2b_1)c_2=a_1y-a_2x.
\]
\end{derivationbox}

\begin{theorembox}[Coordinates in a planar basis]
If
\[
\mathcal B=(\mathbf b_1,\mathbf b_2),
\qquad
\mathbf b_1=(a_1,a_2),
\qquad
\mathbf b_2=(b_1,b_2)
\]
is a basis and $\mathbf u=(x,y)$, then
\[
\boxed{
 c_1=\frac{xb_2-yb_1}{a_1b_2-a_2b_1},
 \qquad
 c_2=\frac{a_1y-a_2x}{a_1b_2-a_2b_1}
}.
\]
\end{theorembox}

\begin{interpretationbox}[Why the denominator matters]
The denominator
\[
a_1b_2-a_2b_1
\]
is nonzero precisely because the basis vectors are not parallel. If it vanished,
the two chosen directions would collapse onto one line and the coordinates would
not be uniquely determined.
\end{interpretationbox}
